In every two-colouring of the positive squares, must every sufficiently large integer be a monochromatic sum of distinct squares? Yes. We use Conlon--Fox--Pham's polynomial Ramsey-completeness theorem in its general form, and verify the ordinary-completeness hypothesis constructively for squares.
The external theorem, Theorem 1.2 of their paper, says: if the sequence of values of a real polynomial $P$ of degree $d$ is complete (every sufficiently large integer is a sum of distinct terms), then for every integer $r\ge2$ it has an $r$-Ramsey-complete subsequence $B$ with
\[
 |B\cap[1,X]|\le C(d)r(\log X)^2
\]
for sufficiently large $X$, with the harmless small-$X$ convention absorbed into the bound. Ramsey completeness means that every $r$-colouring permits monochromatic representations of all sufficiently large integers. The proof of this deep polynomial theorem is external; its completeness hypothesis and its deduction for the present problem are handled below.
